\documentclass[a4paper,10pt]{article}

\usepackage[normalem]{ulem}
\usepackage{contour}
\renewcommand{\ULdepth}{3.4pt}

\usepackage{latexsym}
\usepackage[empty]{fullpage}
\usepackage{titlesec}
\usepackage{marvosym}
\usepackage[usenames,dvipsnames]{xcolor}
\usepackage{verbatim}
\usepackage{enumitem}
\usepackage[hidelinks]{hyperref}
\usepackage{fancyhdr}
\usepackage[polish]{babel}
\usepackage{tabularx}
\usepackage[utf8]{inputenc}
\usepackage[T1]{fontenc}
\usepackage{ragged2e}

\addtolength{\oddsidemargin}{-0.5in}
\addtolength{\evensidemargin}{-0.5in}
\addtolength{\textwidth}{1in}
\addtolength{\topmargin}{-0.75in}
\addtolength{\textheight}{1.5in}

\urlstyle{same}
\raggedbottom
\raggedright
\setlength{\tabcolsep}{0in}

\definecolor{TechBlue}{HTML}{1A4B84}
\definecolor{LightGray}{HTML}{7A7A7A}

\titleformat{\section}{
  \vspace{-2pt}\scshape\raggedright\large\color{TechBlue}
}{}{0em}{}[\color{TechBlue}\titlerule \vspace{2pt}]

\newcommand{\resumeItem}[1]{
  \item\small{
    {#1 \vspace{-1pt}}
  }
}
\newcommand{\resumeSubheading}[4]{
  \vspace{0pt}\item
    \begin{tabular*}{0.97\textwidth}[t]{l@{\extracolsep{\fill}}r}
      \textbf{#1} & #2 \\
      \textit{\small#3} & \textit{\small #4} \\
    \end{tabular*}\vspace{-2pt}
}
\newcommand{\resumeProjectHeading}[2]{
    \vspace{0pt}\item
    \begin{tabular*}{0.97\textwidth}{l@{\extracolsep{\fill}}r}
      \small\textbf{#1} & #2 \\
    \end{tabular*}\vspace{-3pt}
}

\newcommand{\resumeSubHeadingListStart}{\begin{itemize}[leftmargin=0.15in, label={}, topsep=0pt]}
\newcommand{\resumeSubHeadingListEnd}{\end{itemize}\vspace{0pt}}
\newcommand{\resumeItemListStart}{\begin{itemize}[itemsep=1pt, parsep=0pt, partopsep=0pt, topsep=0pt]}
\newcommand{\resumeItemListEnd}{\end{itemize}\vspace{0pt}}

\begin{document}

\begin{center}
    \textbf{\Huge \scshape Jakub Jurkian} \\ \vspace{3pt}
    {\large \color{TechBlue} Java Backend Developer} \\ \vspace{3pt}
    Gdańsk, Polska \\ \vspace{2pt}
    \small
    \href{mailto:kuba.jur03@gmail.com}{\uline{kuba.jur03@gmail.com}} $|$
    +48 724 856 514 $|$
    \href{https://linkedin.com/in/jakub-jurkian}{\uline{LinkedIn}} $|$
    \href{https://github.com/jakub-jurkian}{\uline{Github}}
\end{center}

\section{O mnie}
{\justifying \small Student informatyki na Uniwersytecie Gdańskim, tworzący systemy backendowe w technologiach \textbf{Java} i \textbf{Spring Boot}. Przykładam dużą wagę do dobrze przetestowanego kodu, który można bez obaw przekazać innemu programiście: kluczową logikę piszę w podejściu test-first, a decyzje architektoniczne dokumentuję w ADR-ach. Szukam stażu jako Java Backend Developer lub Software Engineer - dostępny \textbf{od zaraz} w niepełnym wymiarze - lub pracy jako Junior (pełen etat) od \textbf{stycznia 2027 r.}\par}

\section{Umiejętności techniczne}
 \begin{itemize}[leftmargin=0.15in, label={}, topsep=0pt]
    \small{\item{
     \textbf{Języki}{: \textbf{Java 21+}, SQL, Python, TypeScript} \\
     \textbf{Frameworki i testowanie}{: \textbf{Spring Boot 3+}, Spring Security, Hibernate/JPA, Testcontainers, JUnit 5, Mockito} \\
     \textbf{Infrastruktura i bazy danych}{: \textbf{PostgreSQL}, Redis, Liquibase, Docker, GitHub Actions CI/CD, MongoDB} \\
     \textbf{Architektura i metodyki}{: RESTful API, TDD, Architecture Decision Records (ADR), Git, praca oparta na GitHub Issues i PR-ach} \\
     \textbf{Dodatkowo eksploruję}{: React, Node.js/Express}
    }}
 \end{itemize}

\section{Projekty}
    \resumeSubHeadingListStart
      \resumeProjectHeading
          {Fleet Management API (VeloCity) $|$ \emph{Java 25, Spring Boot 4, PostgreSQL} $|$ \href{https://github.com/jakub-jurkian/velocity-api}{\underline{GitHub}} $|$ \href{https://velocityfleet.dev}{\underline{Demo}}}
          {2026}
          \resumeItemListStart
        \resumeItem{Zbudowałem REST API w Spring Boot (z osobnym klientem React/TypeScript) dla wypożyczalni rowerów elektrycznych, zabezpieczone bezstanowym łańcuchem filtrów JWT i blacklistą tokenów w Redis przy wylogowaniu.}
        \resumeItem{Wyeliminowałem double-booking na poziomie bazy danych dzięki PostgreSQL exclusion constraint (GiST na zakresach dat) mapowanemu na HTTP 409; potwierdza to test współbieżności (Testcontainers).}
        \resumeItem{Cykl życia rezerwacji zamodelowałem jako maszynę stanów w warstwie domeny z optimistic locking (@Version) chroniącym przed utratą aktualizacji; schematem bazy zarządzają migracje Liquibase.}
        \resumeItem{Silnik wyceny (BigDecimal) napisałem w podejściu TDD. Ponad 40 testów jednostkowych i integracyjnych działa w GitHub Actions, a osobny workflow wdraża obraz Dockera na maszynę wirtualną w Oracle Cloud.}
          \resumeItemListEnd

      \resumeProjectHeading
          {LLM Evaluation API (MockBean) $|$ \emph{Java 21, Spring Boot 3, LangChain4j, pgvector} $|$ \href{https://github.com/jakub-jurkian/mockbean}{\underline{GitHub}}}
          {2026}
          \resumeItemListStart
            \resumeItem{Zbudowałem API w Spring Boot oceniające odpowiedzi z rozmów technicznych (LLM-as-a-judge) w trybie RAG (pgvector) lub czystego promptu; model (Claude, Gemini, lokalna Llama 3.1) wybiera profil Springa.}
            \resumeItem{Przeprowadziłem benchmark 3 modeli $\times$ 2 trybów na 40 oznaczonych odpowiedziach z przedziałami ufności bootstrap (Python): Claude osiągnął 92,5\% trafności, a RAG nie poprawił wyników żadnego z modeli.}
            \resumeItem{Wdrożyłem migracje schematu bazy danych we Flyway (zamiast Hibernate auto-DDL). System zapisuje każdą ocenę w celu późniejszej analizy spójności modeli.}
          \resumeItemListEnd

      \resumeProjectHeading
          {E-Commerce Marketplace API (Grailkits) $|$ \emph{Express.js, Keycloak, PostgreSQL/MongoDB} $|$ \href{https://github.com/jakub-jurkian/grailkits}{\underline{GitHub}}}
          {2026}
          \resumeItemListStart
            \resumeItem{Podzieliłem backend platformy e-commerce na cztery niezależne serwisy (API Gateway, Catalog, Order, Review) i osobno skonteneryzowałem każdy z nich.}
            \resumeItem{Zabezpieczyłem architekturę przy użyciu Keycloak (OAuth 2.0 + PKCE). API Gateway odpowiada za walidację tokenów JWT oraz rate-limiting oparty na Redis, zanim ruch trafi do usług docelowych.}
            \resumeItem{Wykorzystałem bazę PostgreSQL do obsługi zamówień oraz MongoDB do przechowywania szczegółów produktów i opinii (elastyczny schemat danych).}
          \resumeItemListEnd
    \resumeSubHeadingListEnd

\section{Edukacja}
  \resumeSubHeadingListStart
    \resumeSubheading
      {Uniwersytet Gdański}{Gdańsk, Polska}
      {Informatyka, studia stacjonarne (I stopnia)}{2024 - obecnie}
      \resumeItemListStart
        \resumeItem{\textbf{Wybrane przedmioty:} Algorytmy i struktury danych, Aplikacje przemysłowe, Bazy danych, Programowanie współbieżne, Bezpieczeństwo aplikacji webowych, Architektura mikroserwisów.}
        \resumeItem{\vspace{4pt}\textbf{Obowiązkowe praktyki zawodowe:} Szósty semestr (od stycznia 2027 r.) jest w całości przeznaczony na praktyki zawodowe (720 godzin) bez zajęć dydaktycznych, co gwarantuje moją pełną dyspozycyjność.}
      \resumeItemListEnd
    \resumeSubheading
      {Zespół Szkół im. Macieja Rataja}{Reszel, Polska}
      {Technik informatyk (specjalizacja: administracja i programowanie)}{2019 - 2023}
  \resumeSubHeadingListEnd

\section{Języki}
 \begin{itemize}[leftmargin=0.15in, label={}, topsep=0pt]
    \small{\item{
     \textbf{Angielski}{: B2+ -- codzienny język pracy z dokumentacją techniczną, kodem i komunikacją.} \\
     \textbf{Polski}{: ojczysty.}
    }}
 \end{itemize}

\vfill
\begin{center}
  \tiny \color{LightGray} Wyrażam zgodę na przetwarzanie moich danych osobowych dla potrzeb niezbędnych do realizacji procesu rekrutacji zgodnie z Rozporządzeniem Parlamentu Europejskiego i Rady (UE) 2016/679 z dnia 27 kwietnia 2016 r. w sprawie ochrony osób fizycznych w związku z przetwarzaniem danych osobowych i w sprawie swobodnego przepływu takich danych oraz uchylenia dyrektywy 95/46/WE (RODO).
\end{center}

\end{document}
